La voz humana es una señal biológica compleja resultante de la interacción dinámica entre
la aducción y vibración de los pliegues vocales y la emisión de flujo de aire pulmonar a través
de estructuras resonantes (\cite{lee2021}). Como característica biométrica, posee tanto rasgos
físicos (fisiológicos) como conductuales que la hacen única para cada individuo. En el ámbito
de la seguridad, el reconocimiento de voz se divide principalmente en dos tareas: la
identificación de hablante, que determina la identidad en una comparación de uno a muchos, y
la verificación de hablante, que confirma una identidad declarada mediante un mapeo uno a
uno (\cite{shivani}). Estos procesos pueden ser dependientes del texto o
independientes del mismo, siendo estos últimos más versátiles pero complejos de procesar en
entornos prácticos (\cite{costantini}).

El procesamiento de voz en sistemas embebidos requiere transformar la señal de audio
cruda en representaciones paramétricas compactas que reduzcan la dimensionalidad sin perder
información discriminativa. Las técnicas más prevalentes incluyen a los Coeficientes
Cepstrales en la Escala de Mel (MFCC), que representan el estándar al aproximar la respuesta
no lineal del sistema auditivo humano mediante filtros espaciados logarítmicamente. Son
altamente efectivos para modelar la envolvente espectral y el tono. Por otro lado, también
existen los Coeficientes de Predicción Lineal (LPCC), que se basan en el sistema articulatorio
humano (tracto vocal) y modelan la señal como una combinación lineal de muestras pasadas.
Son computacionalmente eficientes y útiles en tareas de frase fija. También, se puede contar
con características de dominio temporal (ZCR y STE), donde la tasa de cruce por cero (ZCR)
y la energía de tiempo corto (STE) son descriptores de baja complejidad computacional
ideales para dispositivos con recursos extremadamente limitados, como los
microcontroladores PIC (\cite{shili}).

La implementación de modelos de aprendizaje automático directamente en
microcontroladores (Edge AI) permite reducir la latencia, mejorar la privacidad y eliminar la
dependencia de la conectividad en la nube (\cite{roharia}). En cuanto a los tipos de arquitectura, las
redes neuronales convolucionales (CNN) son altamente efectivas al tratar los espectrogramas
de voz como imágenes para identificar patrones locales (\cite{costantini}). Por otro lado, los 
Perceptrones Multicapa (MLP) o redes feed-forward (FNN) ofrecen
un equilibrio óptimo entre precisión y ligereza para tareas de clasificación binaria en hardware
restringido. En cuanto a la optimización, se emplea el método llamado Cuantización. El
mismo es empleado para adaptar modelos complejos a la memoria SRAM y Flash limitada de
un microcontrolador, que convierten pesos de punto flotante de 32 bits a enteros de 8 bits
(int8), reduciendo el tamaño del modelo hasta cuatro veces con una pérdida mínima de
precisión (\cite{evans}).

El rendimiento de los sistemas de biometría vocal es sensible a variables ambientales y
biológicas como:
\begin{itemize}
    \item  Relación Señal-Ruido (SNR): la calidad del micrófono MEMS es crítica; micrófonos con
alto SNR permiten capturar sutilezas de la voz incluso a niveles bajos de presión sonora
(SPL), mientras que un bajo SNR degrada significativamente la tasa de error de palabra
(WER).
    \item Variabilidad biológica y envejecimiento: el envejecimiento fisiológico altera las
estructuras laríngeas, modificando la frecuencia fundamental (F0), el jitter y el shimmer.
Estadísticas de orden superior, como la asimetría (skewness) y la curtosis, han demostrado ser
predictores clave para distinguir entre voces jóvenes y ancianas.
    \item Vulnerabilidades de Seguridad: Los sistemas son susceptibles a ataques de replaying y
clonación de voz mediante IA generativa. Como contramedida, se implementan técnicas de
detección de vitalidad (liveness detection) y autenticación multifactor (MFA) combinando voz
con otros rasgos biométricos.
\end{itemize}

La eficacia de un sistema biométrico se cuantifica mediante la tasa de falsa aceptación
(FAR), la tasa de falso rechazo (FRR) y la tasa de igual error (EER) (\cite{brydinskyi}).
El EER representa el punto donde FAR y FRR son iguales; un valor menor indica una mayor
precisión global del sistema. En aplicaciones de comercio electrónico o control de acceso,
estas métricas deben balancearse según la criticidad de la seguridad frente a la conveniencia
del usuario.

\subsubsection{Modelos de identificación del hablante}
\textcolor{red}{En este apartado se detallará acerca de la arquitectura de los distintos tipos de 
modelos que se encuentran en el mercado.}

\subsubsection{Embeddings}
\textcolor{red}{En este apartado se detallará acerca de la herramienta fundamental que es 
para esta investigación la utilización de embeddings.}

\subsubsection{Espectrograma de Mel y MFCC}
\textcolor{red}{En este apartado se detallará acerca del concepto fundamental para la extracción de 
información acústica de la señal que ingresa al sistema.}